\section{Experimental Setup}\label{sec:5-chapter-experimental-setup}
% In this section, we describe the experimental setup for our three research questions.
\subsection[Current PLS evaluation standards]{Current PLS evaluation standards \hyperref[rq1]{\rqone}}

We aim to conduct a thorough literature survey of the standard practices in readability evaluation for PLS. We collect papers\footnote{On May 7th, 2025} from the ACL Anthology\footnote{\url{https://aclanthology.org/}} that mention one of the following key phrases: ``plain language summarization,'' ``readable summaries,'' or ``lay summarization.'' We exclude papers published for a shared task from annotation and assume the participants use the metrics designated by the shared task organizers. Our goal is to understand the decisions made by researchers, and including shared task papers in this survey would over-represent the decisions made by the task organizers. We report the evaluation methods used by the shared tasks and the number of participants to represent the impact of the evaluation choices. We identify 55 papers that match our criteria. We annotate the papers for relevance to PLS, the type of publication (Main conference, Findings, or Workshop), and which readability evaluation metrics are used. We exclude papers from the survey not relevant to PLS, resulting in 18 relevant papers from the years 2013 to 2025. The most common reasons for relevance exclusion include using ``readable'' in a different word sense (e.g. ``human readable'' vs ``machine readable'') or just citing a PLS paper. We report the number of papers that use each  metric.

\subsection[Comparing traditional readability metrics to human judgments]{Comparing traditional readability metrics to human judgments \hyperref[rq2]{\rqtwo}}

\paragraph{Human Annotated Data.} 
To measure the correlation between readability metrics with human judgments, we use the dataset collected by ~\textcite{August2024KnowYA}. This dataset contains 60 summaries of 10 scientific papers in a variety of domains. Each paper has both expert written and machine written summaries (generated using GPT-3.) The summaries are annotated on a scale of 1 to 5 for the annotator's reading ease of the article. 1 indicates a very poor reading ease, while 5 indicates a very high reading ease. For each summary, we take the average of the annotators' scores to calculate the correlations with readability evaluations as described below. \textcite{August2024KnowYA} originally collected this dataset to better understand human preferences in scientific summarization. In this chapter, we extend their work by applying their findings to summarization evaluation metrics. To the best of our knowledge, this is the only available dataset of human judgments for PLS. \Cref{appendix:dataset_details} contains additional dataset details. 


\paragraph{Traditional readability metrics.} We consider ``traditional'' readability metrics to be those most commonly used in PLS literature. These metrics are well-established, and have been used in past work as judges of readability~\cite{chandrasekaran2020overview}. This term excludes LM-based evaluations, discussed in~\Cref{sec:rq-three-setup}.
We consider 8 readability metrics: Flesh-Kincaid Grade Level (FKGL)~\cite{flesch1952simplification}, Flesch Reading Ease (FRE)~\cite{Flesch1943MarksOR}, Dale Chall Readability Score (DCRS)~\cite{dale1948formula}, Automated Readability Index (ARI)~\cite{smith1967automated}, Coleman Liau Index (CLI)~\cite{Coleman1975ACR}, Gunning Fog Index (GFI)~\cite{Gunning1968TheTO}, Spache~\cite{spache1953new} and Linsear Write (LW)~\cite{o1966gobbledygook}.
% \daniel{if these have citations, cite here}. 
All of the metrics, except for DCRS and Spache, use lexical features such as number of syllables or length of sentences to measure readability. DCRS and Spache use word familiarity to measure readability, assuming that more common words are easier to read~\cite{dale1948formula,o1966gobbledygook}.%\daniel{either you need to define what familiarity is or you need to cite; otherwise this sentence is incomplete}. 
\footnote{We use the \href{https://github.com/cdimascio/py-readability-metrics}{py-readability-metrics} package to calculate the readability scores.} %All of the metrics listed above, except LW and FRE, provide a lower score for higher readability, while the human judgments provide a higher score for higher readability. 



\paragraph{Quantifying alignment between traditional metrics and humans.} We report the Pearson and Kendall-Tau correlation of each metric listed above with the human judgments collected by~\textcite{August2024KnowYA}.
Except for LW and FRE, all metrics provide a lower score for higher readability, while the human judgments provide a higher score for higher readability. To calculate correlations, we multiply the scores by $-1$ (except for LW and FRE), so that text rated as more readable by traditional metrics will be positively correlated with human judgments.

\subsection[LMs as evaluators of readability]{LMs as evaluators of readability \hyperref[rq3]{\rqthree}}\label{sec:rq-three-setup}

We experiment with the following 5 LMs as evaluators of readability: Mistral 7B~\cite{jiang2023mistral7b}, Mixtral 7B~\cite{Jiang2024MixtralOE}, Gemma 7B~\cite{gemma_2024}, Llama 3.1 8B, and Llama 3.3 70B~\cite{grattafiori2024llama3}. We experiment with 3 prompts and report the prompts in \cref{appendix:prompt}. We report the Pearson and Kendall-Tau correlations of the scores provided by each LM with the human judgments.


\subsection[Analysis of summarization datasets]{Analysis of summarization datasets \hyperref[rq4]{\rqfour}}
% \daniel{Consider adding paragraph heading to the next two paragraphs so that it's clear what each is about.}
To test the ability of our results to generalize to datasets outside of the one collected by \textcite{August2024KnowYA}, we include datasets with intended audiences more specific than ``general'' - experts and kids. We expect expert-targeted datasets to be given low readability scores and kid-targeted datasets to have high readability scores.

\paragraph{Expert targeted datasets.} We include 3 expert-targeted datasets: arXiv, PubMed~\cite{cohan2018discourse} and SciTLDR~\cite{cachola2020tldr}.
arXiv and PubMed are collections of abstracts in the Computer Science and Biomedical domains, respectively~\cite{cohan2018discourse}. SciTLDR is a collection of short, expert-targeted, one sentence summaries of Computer Science papers. We expect our methods to provide low readability scores. Additionally, the comparison of SciTLDR to arXiv and PubMed allows us to test if the scores are length dependent. 

\paragraph{Kid-targeted dataset.} The Science Journal for Kids (SJK) dataset is a collection of summaries of scientific papers in a variety of domains, intended for kids~\cite{BioLaySumm_2024}. 
% We noticed the original SJK dataset published by~\textcite{BioLaySumm_2024} contained a number of grammatical mistakes, likely a result of PDF processing errors. We recollected the dataset by accessing the summaries directly on the SJK website\footnote{\url{https://www.sciencejournalforkids.org/}} to ensure that the summaries are grammatically correct. 
Given that this dataset is targeted to kids, we expect it would receive high readability scores.


\paragraph{General audience datasets.} In addition to the datasets listed above, we evaluate 6 popular datasets intended for PLS: CDSR~\cite{Guo2020AutomatedLL}, PLOS~\cite{goldsack2022making}, eLife~\cite{goldsack2022making}, Eureka~\cite{Zaman2020HTSSAN},  CELLS~\cite{Guo2022RetrievalAO}, and SciNews~\cite{pu-etal-2024-scinews}. These datasets are intended for a general audience. CDSR, PLOS, and CELLS are written by journal editors or experts. eLife Sciences gives paper authors the option to write ``eLife digests,'' with the goal of ``cutting jargon and putting research in context.''
\footnote{\url{https://elifesciences.org/digests}} 
The Eureka dataset was collected from EurekaAlert, which hosts press releases about research for scientific journalists.
% \footnote{\url{https://www.eurekalert.org/}} 
Finally, SciNews is a collection of scientific news reports, written by science reporters.


\Cref{tab:dataset-comp} contains a comparison of the summarization datasets analyzed in this chapter. We use the test split of each dataset for our analysis and we report the intended audience, domain, number of documents in the test set, and average number of white-space delineated tokens.
In total, we analyze 10 popular scientific summarization datasets. 


\begin{table}[t]
\setlength{\tabcolsep}{3pt}
\centering \footnotesize
\begin{tabular}{r|lll l}
\textbf{Dataset} & \textbf{Audience} & \textbf{Domain} & \textbf{\# Docs} & \textbf{\# Tokens} \\ \hline
% \multicolumn{4}{r}{\textit{Readability Known Datasets}} \\ \hdashline[0.5pt/3pt]
arXiv & Experts & CS &6440 & 163 \\
PubMed & Experts & Medicine &6658 & 205 \\
SciTLDR & Experts & CS &618 & 19 \\
SJK & Kids & Varied & 284  & 142 \\ \hdashline[0.5pt/3pt]
% \multicolumn{4}{r}{\textit{PLS Datasets}} \\ \hdashline[0.5pt/3pt]
CDSR & General & Healthcare &284 & 221 \\
PLOS & General & Biomed &1376 & 195 \\
eLife & General & Biomed &241 & 383 \\
Eureka & Journalists & Varied &1010 & 662 \\
CELLS & General & Biomed &6311 & 162 \\
SciNews & General &  Varied & 4188 & 615 
\end{tabular}
\caption[Comparison of the datasets analyzed in this chapter.]{Comparison of the datasets analyzed in this chapter. The first 4 are datasets with a specific target audience. The following 6 datasets are commonly used in PLS literature. We report the number of documents (\# Docs) in the test set as well as the average number of tokens (\# Tokens).
% \daniel{FYI: change D and T to Docs/Tokens since you haves room.}
}\label{tab:dataset-comp}
\end{table}
